% ============================================================================
%  Saussure in HLLSet Algebra --- Part III: Diachronic Linguistics
%  A computational exegesis of Ferdinand de Saussure's
%  Course in General Linguistics, Part Three, in the Boolean lattice
%  and GF(2) ring of HLLSet Algebra
% ----------------------------------------------------------------------------
%  Series:  Saussure in HLLSet Algebra (reverse direction of influence)
%  Build:   pdflatex saussure_hllset_part3_diachronic_linguistics.tex
%           (run twice)
%  Notation follows Part I:
%    saussure_hllset_part1_general_principles.tex
%  and Part II:
%    saussure_hllset_part2_synchronic_linguistics.tex
%  and the EWM PUB series.
% ============================================================================
\documentclass[11pt,a4paper]{article}

\usepackage[utf8]{inputenc}
\usepackage[T1]{fontenc}
\usepackage{amsmath,amssymb,amsthm}
\usepackage{amscd}
\usepackage{booktabs}
\usepackage{array}
\usepackage{placeins}
\usepackage{tikz}
\usetikzlibrary{arrows.meta,positioning}
\usepackage{hyperref}
\usepackage[margin=1in]{geometry}
\usepackage{enumitem}
\usepackage{xcolor}
\definecolor{hlblue}{RGB}{0,80,160}
\hypersetup{colorlinks=true,linkcolor=hlblue,urlcolor=hlblue,citecolor=hlblue}

% ---- theorems (aligned with the PUB series, Part I, and Part II) -----------
\newtheorem{principle}{Principle}[section]
\newtheorem{theorem}[principle]{Theorem}
\newtheorem{definition}[principle]{Definition}
\newtheorem{proposition}[principle]{Proposition}
\newtheorem{corollary}[principle]{Corollary}
\newtheorem{remark}[principle]{Remark}
\theoremstyle{remark}
\newtheorem{observation}[principle]{Observation}
\newtheorem{reading}[principle]{Reading}

% ---- shorthands (aligned with Parts I--II and the PUB series) ---------------
\newcommand{\Plane}{\ensuremath{\mathbb{B}}}            % bit plane
\newcommand{\Lat}{\ensuremath{\mathcal{B}}}              % HLLSet (Boolean) lattice
\newcommand{\Tok}{\ensuremath{\mathcal{T}}}              % token space
\newcommand{\HLL}{\ensuremath{H}}                        % an HLLSet
\newcommand{\xor}{\ensuremath{\Delta}}                   % ring addition
\newcommand{\BSS}{\ensuremath{\mathsf{BSS}}}             % BSS measure
\newcommand{\res}{\ensuremath{\mathrm{res}}}             % residual
\newcommand{\spanof}[1]{\ensuremath{\langle #1\rangle}}  % GF(2) span
\newcommand{\supp}{\ensuremath{\mathrm{supp}}}           % support
\newcommand{\Ctx}{\ensuremath{\mathcal{C}}}              % token context
\newcommand{\Ing}{\ensuremath{\mu}}                      % ingestion morphism
\newcommand{\Mat}{\ensuremath{\rho}}                     % materialization morphism
\newcommand{\Lang}{\ensuremath{\mathcal{L}}}             % langue (the lattice system)
\newcommand{\Par}{\ensuremath{\mathcal{P}}}              % parole (usage)
\newcommand{\Sign}{\ensuremath{\mathcal{S}}}             % the sign relation
\newcommand{\Cov}{\ensuremath{\mathrm{cov}}}             % coverage
\newcommand{\Disp}{\ensuremath{\mathcal{D}}}             % the dispatcher
\newcommand{\Sync}{\ensuremath{\mathsf{Syn}}}            % synchronic fiber
\newcommand{\Dia}{\ensuremath{\mathsf{Dia}}}             % diachronic trajectory
\newcommand{\Gen}{\ensuremath{\gamma}}                   % basis generation
\newcommand{\cover}{\ensuremath{\mathrm{cover}}}         % frame cover
\newcommand{\Val}{\ensuremath{\mathcal{V}}}              % value (differential)
\newcommand{\Synt}{\ensuremath{\mathsf{Synt}}}           % syntagmatic reading
\newcommand{\Ass}{\ensuremath{\mathsf{Ass}}}             % associative reading
% ---- new shorthands for Part III -------------------------------------------
\newcommand{\Dep}{\ensuremath{\mathsf{Dep}}}             % departed
\newcommand{\Ret}{\ensuremath{\mathsf{Ret}}}             % retained
\newcommand{\New}{\ensuremath{\mathsf{New}}}             % new
\newcommand{\Trans}{\ensuremath{\mathsf{Tr}}}            % recorded transition
\newcommand{\Piv}{\ensuremath{\mathsf{Piv}}}             % pivot
\newcommand{\Ana}{\ensuremath{\mathsf{Ana}}}             % analogical fourth
\newcommand{\Change}{\ensuremath{\mathbf{M}}}            % change-of-basis matrix

% ---- title -----------------------------------------------------------------
\title{\textbf{Saussure in HLLSet Algebra}\\[4pt]
\large Part III: Diachronic Linguistics\\[4pt]
\large A Computational Exegesis of the \textit{Course in General Linguistics},\\
Part Three, in the Boolean Lattice and GF(2) Ring of HLLSet Algebra\\[6pt]
\normalsize Third paper in the series \textit{Saussure in HLLSet Algebra}}

\author{Alex Mylnikov \and Alexandr Mylnikov Jr. \and Deependra Kumar}
\date{\today}

\begin{document}
\maketitle

\begin{abstract}
Part III of the series releases the exclusion that Part II imposed: time.
Part Three of the \textit{Course}, \emph{Diachronic Linguistics}, is re-founded
as the theory of the \textbf{pivot}. The center of the exegesis is the
generation change of the basis frame. A phonetic change is a pivot of the
ring---an \emph{extension} (a new opposition is born) and a \emph{rotation}
(old oppositions are re-pivoted)---and its ``absolute regularity'' is the
determinism of the Gaussian step, while its ``blindness'' is the fact that the
pivot reads only bit positions, never values. Analogy, the counterpart of
sound change, is the \emph{fourth proportional of GF(2)}: given three nodes
$A,B,C$ of the span, the analogical fourth is
$D = A \xor B \xor C$, the unique node with $A \xor B = C \xor D$; it creates
inside the span and never moves the basis. Folk etymology is a \emph{false}
fourth proportional, whose error the algebra measures exactly as
$C \xor C'$. Agglutination is the promotion of a \emph{view} to an
\emph{original}: a frequently welded join is re-ingested as a new unit,
and---because the span is closed under $\xor$ but not under $\cup$---
agglutination is the second source of basis growth. Diachronic units,
identities, and realities close the part: a diachronic unit is a path of
values across generations, diachronic identity is IICA persistence, and
diachronic reality is the recorded transition. The paper closes with the
complete dictionary of Part Three and the residual handed to Part~IV:
geographical diversity.
\end{abstract}

\noindent\textbf{Keywords:} Saussure; diachronic linguistics; sound change;
phonetic change; analogy; folk etymology; agglutination; diachronic unit;
diachronic identity; pivot; basis frame; generation; change of basis; GF(2)
ring; Boolean lattice; HLLSet; computational exegesis.

% ============================================================================
\section{Introduction}
\label{sec:intro}
% ============================================================================

\subsection{Where Part III stands in the series}

Part I of this series \cite{part1} re-founded Part One of the
\textit{Course}---the sign, arbitrariness, linearity, immutability and
mutability, and the static/evolutionary distinction---as theorems and
invariants of HLLSet Algebra. Part II \cite{part2} re-founded Part Two,
Synchronic Linguistics, as the differential theory of value: relative to a
basis frame $F=(B_1,\dots,B_k)$, the hard key $c \in \mathrm{GF}(2)^k$ is the
exact compression of a node's value, identity and value coincide on the span,
and ``only differences, without positive terms'' became the completeness and
difference theorems of \cite{part2}. Both papers held one variable fixed: \emph{the
frame}. Part II said so explicitly: the spill was treated as a boundary, not
as a cause, and the residual handed to Part III was named---the pivot itself
\cite{part2}.

Part Three of the \textit{Course}, \emph{Diachronic Linguistics}, is where
Saussure lets time in. Its subject is not relations between coexisting terms
but relations between \emph{successive} terms: sound changes, their
grammatical consequences, analogy, folk etymology, agglutination, and the
diachronic counterparts of units, identities, and realities. Part III of this
series re-founds each of those chapters in the algebra, around a single
executable event: the \textbf{pivot}, the recomputation of the basis frame
that increments the generation $\Gen$ and re-coordinates every cached value.
The thesis of this paper is one sentence: \textbf{diachronic linguistics is
the theory of the pivot and its two species---the blind pivot (sound change)
and the exact non-pivot (analogy)}. Everything else in Part Three is a
corollary of that distinction.

\subsection{Method, unchanged}

The method is that of Parts I and II \cite{part1,part2}: no metaphor,
principles as theorems or invariants, and ambiguity recorded as data. The
translation note of Part I applies unchanged: we cite Roy Harris's
translation \cite{harris1983}, keep the standard pair
\emph{signifier}/\emph{signified} (Harris: \emph{signal}/\emph{signification}),
and retain the Russian edition \cite{saussure1977} for verification.
References to the \textit{Course} are by part, chapter, and section.

\subsection{What Part Three contains, and the roadmap}

Part Three of the \textit{Course} has eight chapters: (I) general
observations on the diachronic perspective; (II) sound changes; (III)
grammatical consequences of phonetic evolution; (IV) analogy; (V) analogy and
evolution; (VI) popular etymology; (VII) agglutination; (VIII) diachronic
units, identities, and realities, with the appended discussion of subjective
and objective analysis folded into Chapter~VIII here.

Section~\ref{sec:prelim} fixes the standing definitions inherited from Parts I
and II and adds the diachronic layer: transitions, pivots, and the
change-of-basis transport. Section~\ref{sec:ch1} re-founds the general
observations (the diachronic perspective as the trajectory). Section
\ref{sec:ch2} is the center: sound change as the pivot, regularity as
determinism, blindness as bit-level reading. Section~\ref{sec:ch3} derives the
grammatical consequences (bonds broken by transport). Section~\ref{sec:ch4}
is the second center: analogy as the fourth proportional of GF(2). Section
\ref{sec:ch5} separates the two forces of evolution. Section~\ref{sec:ch6}
re-founds folk etymology as a false fourth proportional. Section
\ref{sec:ch7} re-founds agglutination as the promotion of a view to an
original. Section~\ref{sec:ch8} closes with diachronic units, identities, and
realities. Sections \ref{sec:dictionary}--\ref{sec:conclusion} collect the
dictionary, state what Part III adds, and hand the geographical residual to
Part IV.

% ============================================================================
\section{Preliminaries: the synchronic inheritance and the diachronic layer}
\label{sec:prelim}
% ============================================================================

We restate, with their labels, the standing definitions of Parts I and II;
proofs and motivation are in \cite{part1,part2}. We then add the diachronic
layer: transitions, pivots, and transport.

\begin{definition}[Plane, lattice, ring]
\label{def:plane}
The \textbf{plane} is $\Plane = \{0,\dots,32767\}$; the \textbf{HLLSet
lattice} is $\Lat = \mathcal{P}(\Plane)$, a Boolean lattice under
$\cup,\cap,\subseteq$ and a Boolean ring under $\xor$ (written $\Delta$) and
$\cap$. Cardinalities are $|X|$; the complement of $X$ in $\Plane$ is
$\overline{X}$.
\end{definition}

\begin{definition}[IICA morphism; ingestion; materialization]
\label{def:iica}
$\Ing : \Tok \to \Lat$ maps each token to its \textbf{bit-set} and is an
\textbf{IICA morphism} if it is idempotent
(deterministic), immutable, and content-addressed; it extends to finite token
sets by $\Ing(X)=\bigcup_{t\in X}\Ing(\{t\})$. \textbf{Materialization}
$\Mat : \Lat \dashrightarrow \Ctx$ is the partial inverse that restores an
ordered token context through the door, reading each set bit through the LUT
fiber that disambiguates it \cite{attention2026}. (The LUT, multi-hashing,
and the global gates $G_1,G_2,G_3$ are introduced in Part I, Section~2.)
\end{definition}

\begin{definition}[BSS; frames; generations; keys; spill]
\label{def:frame}
Against a reference $B_w \subseteq \Plane$, the \textbf{directed soft key} is
$\BSS_w(X) = |X \cap B_w| / |B_w|$. A \textbf{basis frame} is a row-echelon
GF(2) basis $F=(B_1,\dots,B_k)$ with distinct leading bits, spanning a set of
observed nodes, and stamped with a \textbf{generation} $\Gen$ that increments
on recomputation. Relative to $F$, each $X \in \spanof{F}$ has a unique
\textbf{hard key} $c \in \mathrm{GF}(2)^k$ with
\[
  X \;=\; \bigoplus_{i : c_i = 1} B_i,
\]
the \textbf{soft key} $w_i = \BSS_{B_i}(X)$ for $i=1,\dots,k$, and the
\textbf{spill} $s = |X \setminus \cover(F)|$, where
$\cover(F) = \bigcup_i B_i$ \cite{state2026}.
\end{definition}

\begin{definition}[Synchronic fiber; diachronic trajectory]
\label{def:syn-dia}
The \textbf{synchronic fiber} at time $t$ is $\Sync(t) = (\Lat, F(t),
\Gen(t))$; the \textbf{diachronic structure} is the family $\Dia = \{\Sync(t)
: t \in T\}$ with the recorded transitions between consecutive fibers.
\end{definition}

\begin{definition}[The sign]
\label{def:sign}
The \textbf{sign} is the pair $\Sign = (\Ing,\Mat)$; for a node $A \in \Lat$,
the signifier is the bit-set $A$ and the signified is its differential value
plus the materialized context $\Mat(A)$ \cite{part1,part2}.
\end{definition}

\subsection{The diachronic layer}

Part III needs three objects that Parts I and II used but did not isolate:
the transition, the pivot, and the transport.

\begin{definition}[Transition; the D/R/N triad]
\label{def:transition}
Let $H(t{-}1)$ be the committed base of the previous turn and $S(t)$ the state
after the join of turn $t$. The \textbf{recorded transition} is
\[
  \Trans(t) \;=\; \bigl(\,S(t),\; H(t{-}1),\;
  \Dep,\; \Ret,\; \New\,\bigr),
\]
where the three components are the sets
\[
  \Dep \;=\; H(t{-}1) \setminus S(t), \qquad
  \Ret \;=\; H(t{-}1) \cap S(t), \qquad
  \New \;=\; S(t) \setminus H(t{-}1),
\]
the \textbf{departed}, \textbf{retained}, and \textbf{new} parts of the
transition; they partition $H(t{-}1) \cup S(t)$. The \textbf{coverage} of the
transition is $\Cov(t) = |\Ret|/|S(t)|$ when $|S(t)|>0$.
\end{definition}

\begin{remark}[One letter, two roles, now separated]
\label{rem:drn}
Following the EWM series, the letters $D/R/N$ are reserved for the
transition decomposition of Definition~\ref{def:transition}
(departed/retained/new) \cite{state2026}. The reading of a context $X$
against a document $D$---the numbers $N_D=|D\setminus X|$,
$R_D=|D\cap X|$, $\Cov=R_D/|D|$---is called the \emph{coverage reading}
throughout the series, in Parts I and II as here. The two roles share one
letter no longer.
\end{remark}

\begin{definition}[Pivot: extension and rotation]
\label{def:pivot}
Let $F(t{-}1)$ be a frame of generation $\Gen(t{-}1)$, and let $Y$ be the
original HLLSet of turn $t$. Reduce $Y$ against $F(t{-}1)$ by one Gaussian
step, obtaining the \textbf{residual} $\res = \res_{F(t{-}1)}(Y)$.
\begin{itemize}[leftmargin=*,nosep]
  \item If $\res = \varnothing$, the turn is an \textbf{in-span event}:
    $F(t)=F(t{-}1)$ and $\Gen(t)=\Gen(t{-}1)$.
  \item If $\res \neq \varnothing$, the turn is a \textbf{pivot}. Let $p$ be
    the leading bit of $\res$. \textbf{Extension} adds $B_{k+1}=\res$ as a new
    basis vector with pivot $p$; \textbf{rotation} replaces every existing
    basis vector $B_i$ that contains $p$ by $B_i \xor \res$. The new frame
    $F(t)$ is the resulting row-echelon basis with distinct leading bits, and
    $\Gen(t)=\Gen(t{-}1)+1$.
\end{itemize}
\end{definition}

\begin{proposition}[The pivot is exact]
\label{prop:pivot}
Let $F(t{-}1)$, $Y$, and $\res$ be as in Definition~\ref{def:pivot}.
\begin{enumerate}[leftmargin=*,nosep]
  \item $\res \neq \varnothing$ if and only if $Y \notin \spanof{F(t{-}1)}$.
  \item After a pivot, $\dim \spanof{F(t)} = \dim \spanof{F(t{-}1)} + 1$; the
    old span is contained in the new span, and every old basis vector keeps
    unique coordinates in $F(t)$.
  \item The pivot map $(F, Y) \mapsto F'$ is deterministic: in a fixed
    coordinate order of the plane, the same basis and the same incoming bits
    always produce the same extension and the same rotation.
  \item Spill $|Y \setminus \cover(F(t{-}1))| > 0$ implies
    $\res \neq \varnothing$, but the converse can fail: a node may lie inside
    the cover yet outside the span (the meet obstruction
    \cite{state2026}). The exact trigger of the pivot is $\res \neq
    \varnothing$, not positive spill.
\end{enumerate}
\end{proposition}

\begin{proof}
(1) is the definition of the residual: the remainder vanishes exactly for
span elements. (2) The leading bit $p$ of $\res$ is, by construction, not a
leading bit of $F(t{-}1)$, so $\res$ is linearly independent of the old
basis; the span grows by exactly one, and monotonicity of the span gives the
containment. (3) The Gaussian step and the rotation rule are deterministic
functions of the bit patterns involved. (4) Every span element is a subset of
$\cover(F)$ (an XOR of subsets is a subset of their union), so a bit outside
the cover forces $Y \notin \spanof{F}$; the converse fails because
$B_i \cap B_j \subseteq \cover(F)$ while $B_i \cap B_j \notin \spanof{F}$ in
general.
\end{proof}

\begin{definition}[Change-of-basis transport]
\label{def:transport}
Let $F(t{-}1)=(B_1,\dots,B_k)$ and $F(t)=(B'_1,\dots,B'_{k'})$ be two frames
of a monotone ring, so $\spanof{F(t{-}1)} \subseteq \spanof{F(t)}$. The
\textbf{change-of-basis matrix} $\Change = \Change(t)$ is the
$k \times k'$ matrix over $\mathrm{GF}(2)$ whose $i$-th row is the coordinate
vector of $B_i$ in $F(t)$. A hard key $c$ at $t{-}1$ \textbf{transports} to
the hard key
\[
  c' \;=\; c \,\Change
\]
at $t$ (row-vector convention), and this transport is exact.
\end{definition}

\begin{proposition}[Exact transport]
\label{prop:transport}
With the notation of Definition~\ref{def:transport}, if
$X = \bigoplus_{i:c_i=1} B_i$, then
$X = \bigoplus_{j:c'_j=1} B'_j$ with $c' = c\Change$. Consequently the
coordinates of a fixed node change with the frame, but the change is a
deterministic linear map recorded in the ring, and every reading is valid
relative to the generation in which it was taken.
\end{proposition}

\begin{proof}
Each old basis vector expands uniquely in the new basis; substituting these
expansions into $X$ and collecting coefficients gives $c' = c\Change$ by
linearity over $\mathrm{GF}(2)$.
\end{proof}

% ============================================================================
\section{Chapter I: General observations}
\label{sec:ch1}
% ============================================================================

Saussure opens Part Three by delimiting the diachronic perspective: it studies
relations between \emph{successive} terms that replace one another in time,
not relations between coexisting terms; its facts are events, not states
\cite[Part III, ch.~I]{harris1983}.

\begin{reading}[The diachronic perspective]
\label{read:ch1}
The relations between coexisting terms are the fiber readings of Part II
(hard keys, soft keys, containment). The relations between successive terms
are the recorded transitions $\Trans(t)$. The diachronic perspective is the
reading of the trajectory $\Dia$ through its transitions, exactly as the
synchronic perspective is the reading of one fiber $\Sync(t)$.
\end{reading}

\begin{observation}[Events are not states]
\label{obs:events}
A synchronic fact is a state: it can be read, and re-read, inside one fiber.
A diachronic fact is an event: it happens once, between two fibers, and is
recorded. The two perspectives of Part I---the axis of simultaneities and the
axis of successions---are not two ways of looking at the same object; they
are two different types of object. The algebra makes this a type distinction:
$\Sync(t)$ is a triple $(\Lat, F(t), \Gen(t))$; $\Trans(t)$ is a five-tuple
of nodes with a direction.
\end{observation}

\begin{proposition}[The diachronic law is the recorded transition]
\label{prop:dia-law}
Each transition $\Trans(t)$ is \emph{imperative}---it happened, and the commit
DAG records it---but \emph{not general}---no identity of the lattice forces
it, and it holds of that transition alone. This is Part I's Theorem on the
two kinds of law \cite{part1}, restated with the transition record made
explicit: the diachronic laws are exactly the recorded transitions, and the
synchronic laws are exactly the identities of the fiber.
\end{proposition}

% ============================================================================
\section{Chapter II: Sound changes as pivots}
\label{sec:ch2}
% ============================================================================

Saussure's chapter on sound changes makes two claims that have puzzled
readers. First, sound changes are \emph{absolutely regular}: a sound change
is not an isolated accident but a law holding of every occurrence of the
sound in question, under stated conditions. Second, sound changes are
\emph{blind}: they alter the material of words without any regard for grammar
or meaning \cite[Part III, ch.~II]{harris1983}.

\begin{reading}[Sound change is a pivot]
\label{read:sound}
A sound change is a pivot of Definition~\ref{def:pivot}: a turn whose
original $Y$ is not in the span forces the frame to recompute. The
\emph{extension} is the birth of a new opposition (a new distinctive
coordinate); the \emph{rotation} is the re-pivoting of the old oppositions
that contain the new pivot bit. The generation $\Gen$ increments, and every
hard key in the span is re-coordinated.
\end{reading}

\begin{observation}[A sound change is local in the plane, global in the frame]
\label{obs:local-global}
Let $X$ be a signifier before the change and $X'$ the same signifier after.
A sound change replaces one sound by another; in the algebra the replaced
sound was mapped by $\Ing$ to one bit $b$, and the new sound is mapped to a
different bit $b' \notin X$. The surrounding bits---everything the previous
state of the HLLSet had set around $b$---are retained, so
\[
  X' \;=\; (X \setminus \{b\}) \cup \{b'\},
  \qquad
  X \xor X' \;=\; \{b, b'\},
  \qquad
  \BSS_X(X') \;=\; \frac{|X|-1}{|X|} \;=\; 1 - \frac{1}{|X|}.
\]
The two versions are therefore almost identical as lattice nodes: their
symmetric difference has cardinality $2$, and their directed similarity is
$1 - 1/|X|$, arbitrarily close to $1$ as the node grows. The minor bit
change survives at the lattice level: the new form stays in the same
neighbourhood of the Boolean lattice, and the retained bits keep it
co-occurring with its old context. Yet these same two bits, when $b'$ lies
outside the span, force a pivot: the generation increments and every hard
key in the fiber is re-coordinated. A sound change is thus local in the
plane---two bits---and global in the frame---all values. It is the chess
move of Part I made literal: one piece moves, the whole position is
re-valued.
\end{observation}

\begin{proposition}[Phonetic change as extension and rotation]
\label{prop:sound}
Let $\res \neq \varnothing$ be the residual of turn $t$. Then the pivot acts
as follows. (i) \emph{Extension:} the span dimension grows by one; the
residual becomes a new basis vector, a new distinctive opposition.
(ii) \emph{Rotation:} exactly the old basis vectors containing the new pivot
bit are replaced by their XOR with the residual; the others are untouched.
The stored nodes are never rewritten---IICA immutability---so a sound change
is a change of the frame around the signs, never of the signs themselves.
\end{proposition}

\begin{proof}
This is the Gaussian step of \cite{state2026}, restated in the language of
Definition~\ref{def:pivot}. The rotation count is the number of old
oppositions that re-index themselves through the new one; it is zero exactly
when no old basis vector contains the pivot bit.
\end{proof}

\begin{proposition}[Regularity and blindness, re-founded]
\label{prop:regularity}
(i) \emph{Absolute regularity.} The pivot map of Proposition~\ref{prop:pivot}
is a function: in a fixed coordinate order, the same basis and the same
incoming bits always produce the same extension and rotation. Saussure's
``absolute regularity'' is therefore determinism of the Gaussian step.
(ii) \emph{Blindness.} The pivot reads only bit positions---the leading bits,
the membership of the pivot bit in old vectors, the XOR rule. It never
consults the LUT, the materialized context, the hard keys, or any value.
A sound change is blind because it is a computation on the plane, not on the
system of values.
\end{proposition}

\begin{remark}[The trigger of sound change is a residual]
\label{rem:trigger}
The exact trigger of a sound change is $\res \neq \varnothing$
(Proposition~\ref{prop:pivot}(4)), stated uniformly in Part I \cite{part1}
and here. A node with positive spill is visibly new (bits no opposition can
express); a node inside the cover but outside the span---the meet
obstruction---is invisibly new, and it pivots just the same. Saussure's
sound changes have the same two faces: some are visible to speakers (a new
sound appears), others are visible only to the system (a distinction
reorganizes without new material).
\end{remark}

\begin{table}[htbp]
\centering
\renewcommand{\arraystretch}{1.25}
\begin{tabular}{@{}lll@{}}
\toprule
Object & Before the pivot ($F(0)$, $\Gen=0$) & After the pivot ($F(1)$, $\Gen=1$) \\
\midrule
Plane & \multicolumn{2}{l}{$\{0,1,2,3\}$} \\
Basis & $B_1=\{0,1\},\ B_2=\{1,2,3\}$ &
      $B_1=\{0,1\},\ B_3=\{0,2\},\ B'_2=\{0,1,3\}$ \\
Incoming original & $Y=\{0,2\}$ & residual $\res=Y=\{0,2\}$, pivot bit $2$ \\
Extension & --- & $B_3 = \res$ \\
Rotation & --- & $B'_2 = B_2 \xor \res = \{0,1,3\}$ \\
Node $X = B_1 \xor B_2 = \{0,2,3\}$ & hard key $(1,1)$ & hard key $(1,1,1)$ \\
Transport & --- & $B_1 \mapsto (1,0,0)$, \ $B_2 \mapsto (0,1,1)$ \\
\bottomrule
\end{tabular}
\caption{A miniature sound change on a four-bit plane. The pivot extends the
basis by the residual and rotates the one old basis vector containing the new
pivot bit; the stored node $X$ is untouched while its hard key changes from
$(1,1)$ to $(1,1,1)$.}
\label{tab:pivot}
\end{table}

\begin{remark}[The cost of a sound change]
The rotation count and rotation mass of \cite{state2026} are the exact
measures of a sound change's cost: how many oppositions had to re-index
themselves through the new one. A change with rotation zero adds a new
distinction without disturbing any old one---the diachronic analogue of a
purely additive sound. A change with large rotation is the diachronic
analogue of a systematic shift: many old distinctions are re-drawn at once.
\end{remark}

% ============================================================================
\section{Chapter III: Grammatical consequences of phonetic evolution}
\label{sec:ch3}
% ============================================================================

Saussure's chapter III traces what sound change does to grammar: it breaks
grammatical bonds, produces alternations, and never creates new grammatical
forms---only destroys or re-arranges \cite[Part III, ch.~III]{harris1983}.

\begin{reading}[Bonds are coordinate coincidences]
\label{read:bonds}
In Part II, the associative series of a node $X$ is the set of nodes sharing
a coordinate with it, and syntagmatic solidarities are coverage guards. A
\emph{grammatical bond} between two nodes is a shared coordinate in the
current frame. The pivot re-coordinates the whole span, so bonds break or
form without any node being rewritten.
\end{reading}

\begin{proposition}[Bonds under transport]
\label{prop:bonds}
Let $X, Y \in \spanof{F(t{-}1)}$ have hard keys $c, d$, and let $\Change$ be
the change-of-basis matrix of the pivot. In the new frame, $X$ and $Y$ share
coordinate $j$ exactly when
\[
  (c \,\Change)_j \;=\; (d \,\Change)_j .
\]
A pivot therefore breaks a bond $(X,Y,i)$ whenever $c_i = d_i$ but
$(c\Change)_j \neq (d\Change)_j$ for every $j$; and it creates a bond
whenever $c_i \neq d_i$ for all $i$ but some transported coordinates agree.
The grammatical consequences of sound change are exactly these breakings and
creations, computed by one matrix multiplication.
\end{proposition}

\begin{proof}
Immediate from Proposition~\ref{prop:transport}: the new coordinates are
$c\Change$ and $d\Change$, and coordinate agreement is equality of the
corresponding entries.
\end{proof}

\begin{corollary}[Alternation is one node, two generation-stamped values]
\label{cor:alternation}
An alternation---two forms of ``the same'' unit in complementary contexts---is
a single stored node whose value changed under a pivot, read against two
generation stamps. The node is one (IICA); its coordinates are two
(transport). This is the sheet of paper of Part II \cite{part2} moving through
time: one cannot change the value without changing the sign's position in the
system, and one cannot change the stored node at all.
\end{corollary}

\begin{remark}[Sound change destroys, never builds]
Saussure insists that phonetic evolution never creates grammatical forms. The
algebra gives the reason: a pivot only re-coordinates existing nodes; it
introduces no new node. Creation is the work of the other force, analogy,
to which we now turn.
\end{remark}

% ============================================================================
\section{Chapter IV: Analogy as the fourth proportional}
\label{sec:ch4}
% ============================================================================

Saussure's chapter IV introduces the second diachronic force. Analogy needs a
model and its regular imitation: it takes \emph{three} terms to create a
\emph{fourth}, by the proportion
\[
  \text{or{\=a}tor} : \text{or{\=a}t{\=o}rem} \;=\;
  \text{hon{\=o}s} : \text{hon{\=o}rem},
\]
and the created form is a grammatical, intelligent construction---unlike sound
change, which is blind \cite[Part III, ch.~IV]{harris1983}.

\begin{reading}[Analogy is the parallelogram law]
\label{read:analogy}
The analogical proportion $A : B \;{::}\; C : D$ is read as the equality of
differences
\[
  A \xor B \;=\; C \xor D .
\]
The proportion is then an equation in the Boolean ring, and it has a unique
solution for the fourth term: $D = A \xor B \xor C$.
\end{reading}

\begin{definition}[Analogical fourth]
\label{def:analogy}
For nodes $A, B, C \in \Lat$, the \textbf{analogical fourth} is
\[
  \Ana(A,B,C) \;=\; A \xor B \xor C .
\]
The proportion $A:B \;{::}\; C:D$ \textbf{holds} exactly when
$A \xor B = C \xor D$, equivalently $D = \Ana(A,B,C)$.
\end{definition}

\begin{theorem}[The analogy theorem]
\label{thm:analogy}
Let $F$ be a frame and $A,B,C \in \spanof{F}$ with hard keys $a,b,c$. Then:
\begin{enumerate}[leftmargin=*,nosep]
  \item \emph{Existence and uniqueness.} $D = \Ana(A,B,C)$ is the unique node
    of $\Lat$ with $A \xor B = C \xor D$.
  \item \emph{Closure.} $D \in \spanof{F}$, with hard key
    $d = a \xor b \xor c$ (coordinatewise XOR).
  \item \emph{Symmetry.} The proportion is its own mirror: also
    $A \xor C = B \xor D$ holds.
  \item \emph{Conservatism.} The frame is untouched:
    $\Gen$ does not increment, and the span does not change. Analogy creates a
    node; it creates no opposition.
\end{enumerate}
\end{theorem}

\begin{proof}
(1) In the group $(\Lat,\xor)$, the equation $A \xor B = C \xor D$ has the
unique solution $D = A \xor B \xor C$ (add $C$ to both sides; in GF(2),
$C \xor C = \varnothing$). (2) The span is a vector subspace over GF(2), so it
is closed under XOR; expanding in the basis gives $d = a \xor b \xor c$.
(3) $A \xor C = A \xor C \xor (B \xor B) = (A \xor B) \xor (B \xor C)$,
while $B \xor D = B \xor A \xor B \xor C = A \xor C$. (4) No pivot is
involved: $D$ is computed from the existing basis, so $F$ and $\Gen$ are
unchanged.
\end{proof}

\begin{proposition}[Analogy and sound change are type-distinct]
\label{prop:two-forces}
Sound change is a \emph{basis event}: the span changes, the generation
increments, stored nodes are untouched, and their values change. Analogy is a
\emph{span event}: the span is untouched, the generation is unchanged, and a
new node appears with a value in the existing system of oppositions.
Saussure's ``analogy is the counterpart of phonetic change'' is, in the
algebra, a type distinction between the only two kinds of diachronic event.
\end{proposition}

\begin{remark}[Why analogy feels intelligent]
Analogy is the one diachronic operation that uses the system of values
itself---the hard keys---as its raw material. A pivot uses only bits; the
analogical fourth uses coordinates. Intelligence, in Saussure's sense, is
exactly this: the operation consults the relations, not the substance.
\end{remark}

% ============================================================================
\section{Chapter V: Analogy and evolution}
\label{sec:ch5}
% ============================================================================

Saussure's chapter V asks how analogy, which is a grammatical and conscious
procedure, can be a factor of \emph{evolution}---a force of change like sound
change \cite[Part III, ch.~V]{harris1983}.

\begin{proposition}[Analogy is conservative in the system, innovative in the lexicon]
\label{prop:analogy-evolution}
Let $D = \Ana(A,B,C)$ be an analogical creation. Then
$\res_F(D) = \varnothing$ and $|D \setminus \cover(F)| = 0$ for the frame in
which $A,B,C$ lie: an analogical creation never forces a pivot, never adds a
bit outside the frame cover, and never changes the system of oppositions. It
is \emph{conservative} in the system---it reuses the community's existing
coordinates---and \emph{innovative} in the lexicon---it names a node that may
never have been uttered before. Once the creation is used in parole and
committed, it enters the record as a new original; but the frame, and with it
the grammar, stays exactly as it was.
\end{proposition}

\begin{proof}
Theorem~\ref{thm:analogy}(2) gives $D \in \spanof{F}$; hence the residual is
empty. Every span element is a subset of the cover, so the spill is zero.
Committing $D$ adds a node to the store, not a direction to the span.
\end{proof}

\begin{proposition}[The two forces decompose the record]
\label{prop:decomposition}
Every diachronic event is of exactly one of two kinds: a \emph{pivot}
(Definition~\ref{def:pivot}), which increments the generation and changes the
system of oppositions; or an \emph{in-span commit}, which records a node
without changing the frame. Sound change is a pivot; analogy is an in-span
commit; agglutination (Section~\ref{sec:ch7}) can be either, depending on
whether the weld escapes the span. The commit DAG labels each commit with its
generation increment, so the two forces are separable in the record itself.
\end{proposition}

\begin{remark}[The speaker's only creative act]
Part I re-founded ``the language changes although the speakers do not change
it'' as a two-layer observation \cite{part1}. Part III now locates the
speaker's contribution precisely. Speakers do not pivot (the pivot is the
system's response to residue); speakers \emph{do} analogize---and even then
they build with the frame the community already shares. The individual
creative act is real, and it is exactly as conservative as Saussure said.
\end{remark}

% ============================================================================
\section{Chapter VI: Popular etymology as a false fourth proportional}
\label{sec:ch6}
% ============================================================================

Saussure's chapter VI treats folk etymology: an opaque form is re-analyzed by
being compared, wrongly, with something familiar. Unlike analogy, it is
isolated and limited in scope, and it rests on an incorrect evaluation of the
old form \cite[Part III, ch.~VI]{harris1983}.

\begin{reading}[Folk etymology is a false analogy]
\label{read:folk}
A folk etymology is an analogy computed from a \emph{mis-identified} term:
the speaker takes the triple $(A,B,C')$ where $C'$ is not the true third term
$C$ but a node confused with it by soft-key resemblance. The result is the
false fourth proportional $D' = \Ana(A,B,C')$.
\end{reading}

\begin{definition}[False fourth proportional]
\label{def:folk}
Let $C$ be the true third term and $C'$ a node mis-identified with $C$. The
\textbf{folk-etymological creation} is
\[
  D' \;=\; \Ana(A,B,C') \;=\; A \xor B \xor C' .
\]
Its \textbf{error} is the symmetric difference between the false and the true
fourth:
\[
  D \xor D' \;=\; (A \xor B \xor C) \xor (A \xor B \xor C')
  \;=\; C \xor C'.
\]
\end{definition}

\begin{proposition}[The algebra measures the incorrect evaluation]
\label{prop:folk}
(i) $D' = D$ if and only if $C' = C$: the folk etymology is exact exactly
when the identification is exact, and then it simply is analogy.
(ii) The error of the creation equals the error of the identification:
$D \xor D' = C \xor C'$. (iii) The false creation carries no provenance
support for the re-analysis: the recipe's sources are the nodes it actually
uses, and the historical relation it claims is not among them.
\end{proposition}

\begin{proof}
(i) and (ii) are immediate from the group law of $\xor$. (iii) A view records
the provenance of the expression actually computed; the claimed relation to
$C$ appears in no source, so the objective analysis (Section~\ref{sec:ch8})
does not confirm it.
\end{proof}

\begin{remark}[Why folk etymology is isolated]
Analogy is systematic because it is computed in the span, with the shared
coordinates: the same proportion holds for every node with the same keys.
Folk etymology is isolated because it is computed from a soft-key confusion
that is local to one form: the mis-identification $C' \approx_{\BSS} C$ need
not transfer to any other triple. Saussure's limitation of scope is, in the
algebra, the non-functoriality of the confusion map.
\end{remark}

% ============================================================================
\section{Chapter VII: Agglutination as the promotion of a view}
\label{sec:ch7}
% ============================================================================

Saussure's chapter VII treats agglutination: two elements that occur side by
side in the chain merge into one unanalyzable unit. It works only within
syntagmas, it is involuntary, and together with analogy it is the main source
of new forms \cite[Part III, ch.~VII]{harris1983}.

\begin{reading}[Agglutination welds a view into an original]
\label{read:agglutination}
In the EWM storage discipline \cite{state2026}, only originals and base
HLLSets are materialized; every compound is a \emph{view}---a recipe plus
provenance, recreated on read. Agglutination is the event that breaks the
discipline on purpose: a compound that the chain keeps producing is promoted
to an \emph{original}---the join is registered as a stored node with its own
content address, and thereafter it behaves as a unit.
\end{reading}

\begin{definition}[Weld]
\label{def:weld}
Let $A, B \in \Lat$ be nodes whose co-occurrence in materialized contexts is
systematically high (the retention reading of the D/R/N triad, accumulated
over windows). The \textbf{weld} of $A$ and $B$ is the registration of
\[
  O \;=\; A \cup B
\]
as a new original: a node that was previously only a view expression
\texttt{union} over sources $A,B$ now has its own stored bits and its own
content address.
\end{definition}

\begin{proposition}[Agglutination can pivot]
\label{prop:agglutination}
Let $A, B \in \spanof{F}$. Then
\[
  A \cup B \;=\; A \xor B \xor (A \cap B),
\]
so the weld $O = A \cup B$ lies in the span if and only if the meet
$A \cap B$ lies in the span. A weld of in-span nodes therefore \emph{forces a
pivot exactly when the meet is not representable}. Agglutination is thus the
second source of basis growth: sound change pivots by residue, agglutination
pivots by meet.
\end{proposition}

\begin{proof}
The identity $A \cup B = A \xor B \xor (A \cap B)$ holds in any Boolean ring.
Since $A \xor B \in \spanof{F}$ (span closed under XOR), membership of the
union in the span is equivalent to membership of the meet; the rest is
Proposition~\ref{prop:pivot}(1).
\end{proof}

\begin{remark}[Three forces, two of which move the basis]
The diachronic forces now have exact types. Sound change: a pivot by
residue, blind, moves the basis. Agglutination: a pivot by meet, syntagmatic,
involuntary, moves the basis when the weld escapes the span, and records an
in-span original otherwise. Analogy: never a pivot, associative, intelligent,
populates the lexicon within the span. Saussure's chapter VII observation
that agglutination, unlike analogy, works only within syntagmas is the type
statement: agglutination is a join (a chain operation), analogy is a XOR (a
frame operation).
\end{remark}

% ============================================================================
\section{Chapter VIII: Diachronic units, identities, realities}
\label{sec:ch8}
% ============================================================================

Saussure closes Part Three with the diachronic counterparts of Part Two's
units, identities, and realities, and the appended distinction between
subjective and objective analysis \cite[Part III, ch.~VIII]{harris1983}.

\begin{reading}[Diachronic unit, identity, reality]
\label{read:dia-unit}
The synchronic unit is a node read in one fiber. The diachronic unit is the
\emph{path} of that node's values across generations. Diachronic identity is
the persistence of the node itself; diachronic reality is the recorded
transition.
\end{reading}

\begin{definition}[Diachronic unit]
\label{def:dia-unit}
A \textbf{diachronic unit} is a pair $(X, [t_1,t_2])$: a stored node $X$
together with the interval of generations over which it is referenced, read
as the value trajectory
\[
  c(t) \;=\; \text{hard key of } X \text{ in } F(t),
  \qquad t_1 \le t \le t_2,
\]
with the transport matrices $c(t{+}1) = c(t)\Change(t{+}1)$ connecting
consecutive values. It is a path in the trajectory $\Dia$, not a node of the
lattice.
\end{definition}

\begin{proposition}[Diachronic identity is IICA persistence]
\label{prop:dia-identity}
A node at $t_1$ and a node at $t_2$ are diachronically identical if and only
if they are the same bit-set,
\[
  X(t_1) \;=\; X(t_2) .
\]
Their values may differ: the hard keys $c(t_1)$ and $c(t_2)$ are equal if and
only if no pivot between $t_1$ and $t_2$ moved them, i.e.\
$c(t_1)\Change = c(t_1)$. Identity is content; value is coordinates. The
distinction of Part II---identity and value coincide within one fiber---now
has its diachronic counterpart: across fibers, identity persists while value
shifts.
\end{proposition}

\begin{proof}
The first assertion is content addressability (IICA): the same bits name the
same node, and different bits name different nodes. The second is
Proposition~\ref{prop:transport}: values change exactly by the transport
matrices of the intervening pivots.
\end{proof}

\begin{proposition}[Subjective and objective analysis]
\label{prop:subjective-objective}
The \textbf{subjective analysis} of a sign is its reading in the current
frame: the soft key and the hard key of $F(t)$---what a speaker of the
current system can see. The \textbf{objective analysis} is the provenance
chain: the view recipe, the commit DAG, and the transport matrices---what the
record shows. A pivot separates the two: subjectively, ``the same word'' has
simply changed its value; objectively, the change is a recorded transition
with a generation stamp. The algebra keeps both, and falsifies neither.
\end{proposition}

\begin{remark}[The residual of Part Three]
Part Three leaves one residual, and it is the same one Saussure leaves: the
\emph{cause} of sound change. The algebra records that a pivot happened, how
many oppositions it rotated, and what it did to every value; it does not, and
cannot, say why the residue arrived---why this community, at this time,
produced this new material. The trigger is parole (the external stream), and
parole is outside the system by construction. The algebra turns Saussure's
``sound changes have no internal cause'' into a type boundary: the pivot is
inside the system; its occasion is not.
\end{remark}

% ============================================================================
\section{The dictionary of Part Three}
\label{sec:dictionary}
% ============================================================================

\begin{table}[htbp]
\centering
\renewcommand{\arraystretch}{1.2}
\begin{tabular}{@{}>{\raggedright\arraybackslash}p{3.4cm}
                    >{\raggedright\arraybackslash}p{4.2cm}
                    >{\raggedright\arraybackslash}p{7.0cm}@{}}
\toprule
Saussure (Part Three) & In the \textit{Course} & In HLLSet Algebra \\
\midrule
diachronic perspective & successive terms, not coexisting & the trajectory
read through the transitions $\Trans(t)$: Reading~\ref{read:ch1} \\
diachronic fact & an event & the recorded transition
$\Trans(t)=(S(t),H(t{-}1),\Dep,\Ret,\New)$: Definition~\ref{def:transition} \\
diachronic law & imperative but not general & the recorded transition itself:
Proposition~\ref{prop:dia-law} \\
sound change & absolutely regular, blind & the pivot (extension + rotation):
Definitions~\ref{def:pivot}; Proposition~\ref{prop:sound} \\
regularity & no exceptions & determinism of the Gaussian step:
Proposition~\ref{prop:regularity}(i) \\
blindness & no regard for grammar/meaning & the pivot reads bits, never
values: Proposition~\ref{prop:regularity}(ii) \\
grammatical consequences & bonds broken, alternations & transport
$c' = c\Change$ breaks/creates coordinate coincidences:
Proposition~\ref{prop:bonds}; alternation = two generation-stamped values:
Corollary~\ref{cor:alternation} \\
analogy & three terms create a fourth & $\Ana(A,B,C)=A\xor B\xor C$, the
fourth proportional: Definition~\ref{def:analogy}, Theorem~\ref{thm:analogy} \\
analogy vs sound change & intelligent vs blind & span event vs basis event:
Proposition~\ref{prop:two-forces} \\
analogy and evolution & conservative and innovative & zero spill, zero
residual, new lexicon node: Proposition~\ref{prop:analogy-evolution}; the
record decomposes into pivots and in-span commits:
Proposition~\ref{prop:decomposition} \\
folk etymology & incorrect evaluation of an old form & false fourth
proportional $D'=\Ana(A,B,C')$ with error $C\xor C'$:
Definition~\ref{def:folk}, Proposition~\ref{prop:folk} \\
agglutination & syntagmatic neighbours merge & promotion of a view to an
original (weld $O=A\cup B$): Definition~\ref{def:weld}; pivots iff the meet
is unrepresentable: Proposition~\ref{prop:agglutination} \\
diachronic unit & not the synchronic unit & a path of values across
generations: Definition~\ref{def:dia-unit} \\
diachronic identity & same essence across time & IICA persistence of the
bit-set: Proposition~\ref{prop:dia-identity} \\
diachronic reality & the event itself & the recorded transition; objective
analysis (provenance + DAG) vs subjective analysis (current frame):
Proposition~\ref{prop:subjective-objective} \\
\bottomrule
\end{tabular}
\caption{The complete dictionary of Part Three of the \textit{Course} into
HLLSet Algebra.}
\label{tab:dictionary}
\end{table}

% ============================================================================
\FloatBarrier
\section{What Part III adds to Parts I and II}
\label{sec:adds}
% ============================================================================

Part I named the generation mechanism; Part II held the frame fixed. Part III
turns the generation mechanism into a datatype. Five additions matter.

\begin{enumerate}[leftmargin=*,nosep]
  \item \textbf{The pivot is now the diachronic unit.} Definition~\ref{def:pivot}
    makes precise what Parts I and II only pointed at: the recomputation of
    the frame, split into extension and rotation, stamped with $\Gen$.
  \item \textbf{Sound change is a theorem, not an image.}
    Propositions~\ref{prop:sound} and \ref{prop:regularity} give the exact
    sense in which a sound change is regular (deterministic) and blind
    (bit-level). The miniature of Table~\ref{tab:pivot} is checkable by hand.
  \item \textbf{Analogy is a theorem of GF(2).} Theorem~\ref{thm:analogy} is
    the parallelogram law: the analogical fourth is $A \xor B \xor C$, unique,
    in-span, generation-preserving. Saussure's proportion stops being an
    observation about speakers and becomes an operation of the ring.
  \item \textbf{The two forces are type-separated, and a third is located.}
    Sound change = basis event; analogy = span event
    (Proposition~\ref{prop:two-forces}); agglutination = view promotion that
    pivots exactly when the meet is unrepresentable
    (Proposition~\ref{prop:agglutination}).
  \item \textbf{The exact pivot trigger.} Proposition~\ref{prop:pivot}(4)
    records the exact trigger of a sound change: $\res \neq \varnothing$.
    Positive spill is the visible sufficient sign; the meet obstruction is
    the invisible case; and the series now states the refined condition
    uniformly in Parts I and III.
\end{enumerate}

% ============================================================================
\section{Conclusion and outlook to Part IV}
\label{sec:conclusion}
% ============================================================================

Part Three of the \textit{Course} has been re-founded as the theory of the
pivot and its two species. Diachronic linguistics is the reading of the
trajectory through recorded transitions; sound change is the blind,
deterministic pivot; its grammatical consequences are transports of
coordinates; analogy is the exact fourth proportional of GF(2); folk
etymology is the same operation with a mis-identified term, and the algebra
measures the mis-identification; agglutination is the promotion of a view to
an original, the syntagmatic source of new oppositions; diachronic units,
identities, and realities are paths, IICA persistence, and recorded
transitions. Saussure's diachronic linguistics, read in this direction, is no
longer a classification of changes: it is the executable theory of a system
whose frame moves.

What Part Three deliberately leaves untouched is \emph{where} the changes
happen. The plane is one, shared by all speakers; nothing has been said about
which communities pivot together, which share a frame, and which do not. Part
Four of the \textit{Course}, Geographical Linguistics, is exactly that
question: the diversity of languages, the complications of geographical
diversity, its causes, and the propagation of linguistic waves. In the
algebra these become the theory of worlds: hash configurations and codebooks
as dialects, the BSS threshold and the LUT as the boundary of a community,
the dispatcher as the speaker who actualizes the shared langue, and the
bridge of \cite{companion2026} as the wave that crosses from one world to
another. Part IV will re-found geographical linguistics as the multi-world
theory of one lattice; Part V will then close the series with retrospective
linguistics, in which reconstruction becomes time travel through old frames.

\section*{Acknowledgments}
This paper was created in collaboration with, and with the assistance of, an
AI system---DeepSeek (deepseek-v4-pro). The pivot reading of sound change,
the fourth-proportional reading of analogy, the view-promotion reading of
agglutination, and several of the formulations emerged in dialogue with the
assistant, whose drafts the authors revised and take full responsibility for
in their final form.

\begin{thebibliography}{99}

\bibitem{saussure1916} F.~de Saussure,
\textit{Cours de linguistique g\'en\'erale}, Payot, Paris, 1916.
(Edited posthumously by C.~Bally and A.~Sechehaye, with A.~Riedlinger.)

\bibitem{harris1983} F.~de Saussure,
\textit{Course in General Linguistics}, translated and annotated by
Roy Harris, Duckworth, London, 1983.

\bibitem{baskin1959} F.~de Saussure,
\textit{Course in General Linguistics}, translated by Wade Baskin,
Philosophical Library, New York, 1959.

\bibitem{saussure1977} F.~de Sossyur,
\textit{Trudy po yazykoznaniyu} [Works in Linguistics], edited by
A.~A.~Kholodovich, translation of the \textit{Cours} by A.~M.~Sukhotin,
Progress, Moscow, 1977.
(Authors' working copy; used for verification of the readings.)

\bibitem{part1} A.~Mylnikov, A.~Mylnikov Jr., and D.~Kumar,
\textit{Saussure in HLLSet Algebra, Part I: General Principles},
SAUSSURE PUB series, 2026.

\bibitem{part2} A.~Mylnikov, A.~Mylnikov Jr., and D.~Kumar,
\textit{Saussure in HLLSet Algebra, Part II: Synchronic Linguistics},
SAUSSURE PUB series, 2026.

\bibitem{ewm2026} A.~Mylnikov, A.~Mylnikov Jr., and D.~Kumar,
\textit{Emerging World Models: A Structural Framework for Self-Organizing
World Representations}, EWM PUB series, 2026.

\bibitem{companion2026} A.~Mylnikov, A.~Mylnikov Jr., and D.~Kumar,
\textit{The EWM Lattice and Its Worlds}, EWM PUB series, 2026.

\bibitem{state2026} A.~Mylnikov, A.~Mylnikov Jr., and D.~Kumar,
\textit{Transformations and Trajectories of a Content-Addressed State
Machine}, EWM PUB series, 2026.

\bibitem{attention2026} A.~Mylnikov, A.~Mylnikov Jr., and D.~Kumar,
\textit{Attention on Lattice Traversal Is All You Need}, EWM PUB series, 2026.

\bibitem{langue2026} A.~Mylnikov, A.~Mylnikov Jr., and D.~Kumar,
\textit{Langue as a Lattice: Closing the Circle from Statistical to
Structural Linguistics through Content-Addressed Algebras}, EWM PUB series,
2026.

\bibitem{birkhoff1940} G.~Birkhoff,
\textit{Lattice Theory}, American Mathematical Society, 1940.

\bibitem{maclane1971} S.~Mac Lane,
\textit{Categories for the Working Mathematician}, Springer, 1971.

\bibitem{hjelmslev1961} L.~Hjelmslev,
\textit{Prolegomena to a Theory of Language}, University of Wisconsin Press,
1961.

\bibitem{harris1954} Z.~S.~Harris,
\textit{Distributional Structure}, Word, 10(2--3):146--162, 1954.

\bibitem{firth1957} J.~R.~Firth,
\textit{A Synopsis of Linguistic Theory, 1930--1955}, Studies in Linguistic
Analysis, 1957.

\bibitem{vaswani2017} A.~Vaswani, N.~Shazeer, N.~Parmar, J.~Uszkoreit,
L.~Jones, A.~N.~Gomez, \L.~Kaiser, and I.~Polosukhin,
\textit{Attention Is All You Need}, NeurIPS, 2017.

\end{thebibliography}

\end{document}
